\documentclass[11pt]{article}

\usepackage[margin=1in]{geometry}
\usepackage{amsmath}
\usepackage{booktabs}
\usepackage{longtable}
\usepackage{graphicx}
\usepackage[dvipsnames]{xcolor}
\usepackage[colorlinks=true, linkcolor=MidnightBlue, citecolor=MidnightBlue, urlcolor=MidnightBlue]{hyperref}
\usepackage{natbib}
\usepackage{enumitem}

\title{TB Programs, Stakeholders, and Elimination Efforts}
\author{}
\date{}

\begin{document}
\maketitle
\tableofcontents
\newpage

% ============================================================
\section{Major Institutions and Stakeholders}
\label{sec:institutions}
% ============================================================

The global TB response is coordinated across a multi-layered network of multilateral bodies, bilateral donors, research institutions, and civil society organizations.

\subsection{Multilateral and Intergovernmental Bodies}

\begin{description}[style=nextline]

  \item[World Health Organization (WHO) -- Global TB Programme]
  The primary normative body for TB policy. WHO publishes the annual Global Tuberculosis Report, sets treatment and prevention guidelines, coordinates the End TB Strategy, and monitors progress toward global targets. Based in Geneva. 
  
  Website: \url{https://www.who.int/teams/global-tuberculosis-programme}

  \item[Stop TB Partnership]
  Established in 2001, the Stop TB Partnership is a network of approximately 2,000 partner organizations (international, governmental, non-governmental, and patient groups). Hosted by UNOPS in Geneva since 2015. It coordinates the Global Plan to End TB 2023--2030, advocates for financing, and convenes seven thematic Working Groups covering drug-resistant TB, new diagnostics, child and adolescent TB, laboratory systems, and community-led responses.
  
  Website: \url{https://www.stoptb.org}

  \item[The Global Fund to Fight AIDS, Tuberculosis and Malaria]
  The largest multilateral financing mechanism for TB, providing approximately 76\% of all international TB funding. The Global Fund channels resources to national TB programs in low- and middle-income countries, with a renewed four-year collaboration framework signed with Stop TB Partnership in November 2024. 
  
  Website: \url{https://www.theglobalfund.org}

  \item[Unitaid]
  A health financing agency hosted by WHO that accelerates access to diagnostics, drugs, and preventive tools. Currently co-funds the DriveDx4TB project (US\$15.9M, 2022--2026) with FIND for point-of-care TB diagnostics in South Africa, India, Kenya, and Indonesia. 
  
  Website: \url{https://unitaid.org}

\end{description}

\subsection{Bilateral Donors and Implementing Agencies}

\begin{description}[style=nextline]

  \item[USAID]
  The U.S. Agency for International Development funds the Tuberculosis Implementation Framework Agreement (TIFA) project, which builds national TB program capacity, and supports FIND for pediatric TB case-finding initiatives in high-burden settings. 
  
  Website: \url{https://www.usaid.gov/global-health/health-areas/tuberculosis}

  \item[PEPFAR (U.S. President's Emergency Plan for AIDS Relief)]
  Though primarily HIV-focused, PEPFAR is a critical TB partner because HIV co-infection is a leading driver of TB mortality. PEPFAR funds integrated TB/HIV services across Sub-Saharan Africa. 
  
  Website: \url{https://www.state.gov/pepfar/}

  \item[CDC (U.S. Centers for Disease Control and Prevention)]
  The CDC's Division of Tuberculosis Elimination supports domestic elimination and provides technical assistance to high-burden countries through the Global TB Program. 
  
  Website: \url{https://www.cdc.gov/tb/}

\end{description}

\subsection{Product Development Partnerships and Research Organizations}

\begin{description}[style=nextline]

  \item[TB Alliance]
  A non-profit product development partnership dedicated to accelerating the development of new TB drug regimens. Manages a clinical portfolio spanning early-phase compounds to Phase 3 trials, including the SPaL and BPaLM regimens described in Section~\ref{sec:treatment}. \url{https://www.tballiance.org}

  \item[FIND (Foundation for Innovative New Diagnostics)]
  A global alliance that connects countries, donors, and industry to accelerate the development and deployment of TB diagnostic tools. Leads the DriveDx4TB project focused on third-generation LAM tests and molecular point-of-care diagnostics. 
  
  Website: \url{https://www.finddx.org}

  \item[IAVI (International AIDS Vaccine Initiative)]
  Supports clinical development of novel TB vaccines including MTBVAC and conducts early-phase immunology research. 
  
  Website: \url{https://www.iavi.org}

  \item[Gates Medical Research Institute (Gates MRI)]
  Executing the Phase 3 clinical trial of the M72/AS01\textsubscript{E} TB vaccine candidate, which reached full enrollment in April 2025 -- approximately 11 months ahead of schedule and involving $\sim$20,000 participants across 54 sites in five countries. 
  
  Website: \url{https://www.gatesmri.org}

  \item[Serum Institute of India / Max Planck Institute]
  Joint developers of VPM1002, a recombinant BCG strain now in Phase 3 development under Serum Institute sponsorship. 
  
  Website: \url{https://newtbvaccines.org/vaccine/vpm1002/}

  \item[Partners In Health (PIH)]
  A global health NGO operating in over 10 high-burden countries. PIH pioneered community health worker-delivered DR-TB care, achieving some of the highest MDR-TB cure rates on record ($\sim$77\%). Currently leads the arcTB project (US\$7.3M, co-funded by Unitaid) to scale shortened all-oral BPaLM-based regimens across seven countries, with technical support from Harvard Medical School and the Institute of Tropical Medicine Antwerp. \\
  Website: \url{https://www.pih.org/programs/tuberculosis}

\end{description}

\subsection{Civil Society and Patient Advocacy}

\begin{description}[style=nextline]

  \item[TB People]
  The global network of people affected by TB. Plays a recognized role in WHO guideline development and Stop TB Partnership governance, advocating for community-centered, rights-based responses and equitable access to treatment. \\
  Website: \url{https://tbpeople.org}

  \item[Treatment Action Group (TAG)]
  An independent, activist research and policy think tank. Monitors and critiques the TB research pipeline, publishes annual TB pipeline reports, and engages with funders and developers to accelerate access to new drugs and diagnostics. \\
  Website: \url{https://www.treatmentactiongroup.org}

  \item[We Are TB]
  A U.S.-based survivor-led advocacy organization providing free peer support groups, patient navigation, and public education for people in TB treatment and their families. Advocates for increased TB funding and reduced stigma by elevating lived-experience testimony to lawmakers and national conferences. \\
  Website: \url{https://www.wearetb.com}

\end{description}

% ============================================================
\section{Ongoing Prevention Programs}
\label{sec:prevention}
% ============================================================

TB prevention operates along two main axes: vaccination to prevent initial infection or disease progression, and TB preventive treatment (TPT) to eliminate latent TB infection (LTBI) before it activates.

\subsection{Vaccination}

\textbf{BCG (Bacillus Calmette--Gu\'{e}rin)} remains the only licensed TB vaccine, introduced in 1921. It provides strong protection against severe childhood forms (miliary and meningeal TB) but inconsistent protection against pulmonary TB in adults. Universal BCG vaccination at birth is included in the WHO Expanded Programme on Immunization and remains part of national immunization schedules in most high-burden countries.

New vaccine candidates in clinical development (see Section~\ref{sec:treatment}) seek to replace or supplement BCG with broader and more durable protection.

\subsection{TB Preventive Treatment (TPT)}

The 2024 WHO consolidated guidelines on TB preventive treatment (second edition) strongly recommend short-course rifamycin-based regimens as the preferred approach over long-course isoniazid monotherapy:

\begin{longtable}{llll}
\toprule
\textbf{Regimen} & \textbf{Drugs} & \textbf{Duration} & \textbf{Recommendation} \\
\midrule
3HP & Isoniazid + rifapentine & 3 months, once weekly & Strong \\
3HR & Isoniazid + rifampin & 3 months, daily & Strong \\
4R  & Rifampin & 4 months, daily & Strong (alternative) \\
1HP & Isoniazid + rifapentine & 1 month, daily & Conditional \\
6H / 9H & Isoniazid & 6--9 months, daily & Alternative when rifamycins unavailable \\
\bottomrule
\caption{WHO-recommended LTBI treatment regimens (2024 edition).}
\label{tab:tpt}
\end{longtable}

Approximately 80\% of TB cases in low-incidence settings arise from reactivation of untreated LTBI. The UN Political Declaration on TB (2023) targets 90\% TPT coverage for eligible populations by 2027, including 30 million household contacts of TB cases and 15 million people living with HIV.

\subsection{Infection Control and Case-Finding}

Upstream prevention efforts include:
\begin{itemize}
  \item Rapid molecular diagnostics (GeneXpert MTB/RIF, truenat) to shorten the diagnostic delay and reduce transmission.
  \item Active case-finding through household contact investigation, systematic screening of high-risk groups (miners, prisoners, healthcare workers), and community mobile screening.
  \item Airborne infection control in healthcare facilities (ventilation, UV germicidal irradiation, respirators).
  \item Social protection interventions to address catastrophic costs, which the End TB Strategy requires to reach zero by 2035.
\end{itemize}

% ============================================================
\section{Treatments: Available and In Development}
\label{sec:treatment}
% ============================================================

\subsection{Current Standard Regimens}

\subsubsection{Drug-Susceptible TB (DS-TB)}

The standard first-line regimen for drug-susceptible pulmonary TB is 2HRZE/4HR: two months of isoniazid (H), rifampin (R), pyrazinamide (Z), and ethambutol (E), followed by four months of isoniazid and rifampin. Treatment success rates exceed 85\% under programmatic conditions when adherence is supported.

\subsubsection{Drug-Resistant TB (DR-TB)}

Multidrug-resistant TB (MDR-TB, resistant to at least isoniazid and rifampin) and rifampin-resistant TB (RR-TB) require longer and more complex treatment. The BPaL regimen (bedaquiline, pretomanid, linezolid) became a landmark advance:
\begin{itemize}
  \item The TB-PRACTECAL trial demonstrated 89\% treatment success at 6 months, compared to 52\% with previous standard-of-care, with a substantially lower adverse-event burden.
  \item BPaLM adds moxifloxacin for patients without fluoroquinolone resistance, further improving outcomes.
  \item WHO now recommends BPaL/BPaLM as the preferred regimen for adults and adolescents aged $\geq$14 years with MDR/RR-TB.
\end{itemize}

\subsection{Vaccines in Advanced Clinical Development}

\begin{description}[style=nextline]

  \item[M72/AS01\textsubscript{E} (GSK/Gates MRI)]
  Subunit vaccine (fusion protein of two \textit{M.~tuberculosis} antigens adjuvanted with AS01\textsubscript{E}). Phase 2b results showed $\sim$50\% efficacy against progression to active TB in IGRA-positive adults. Phase 3 trial (full enrollment April 2025) involves $\sim$20,000 participants in five high-burden countries. If successful, this would be the first new TB vaccine in 100 years. Website: \url{https://www.gatesmri.org}

  \item[MTBVAC (Biofabri/Universit\'{e} de Toulouse/IAVI)]
  Live attenuated vaccine derived from a human clinical \textit{M.~tuberculosis} isolate; the only live TB vaccine candidate in clinical development. Phase 3 pivotal efficacy trial in newborns ongoing in South Africa ($>$4,000 neonates enrolled); Phase 2b in adolescents and adults across South Africa, Kenya, and Tanzania. Phase 2 trials launched in India (Bharat Biotech, 2024--2025). 
  
  Website: \url{https://newtbvaccines.org/vaccine/mtbvac/}

  \item[VPM1002 (Max Planck Institute/Serum Institute of India)]
  Recombinant BCG strain with rearranged \textit{ureC} gene and \textit{listeriolysin} expression designed to improve immune response breadth. Phase 3 development ongoing; intended to replace BCG in neonatal immunization programs. Website: \url{https://newtbvaccines.org/vaccine/vpm1002/}

\end{description}

\subsection{Drug Pipeline}

\begin{description}[style=nextline]

  \item[SPaL (sorfequiline [TBAJ-876] + pretomanid + linezolid)]
  Next-generation regimen under development by TB Alliance. Phase 2 trial (NC-009) demonstrated superior bactericidal activity compared to standard HRZE with a comparable safety profile. Phase 3 entry expected in 2026. Website: \url{https://www.tballiance.org/our-pipeline-portfolio/}

  \item[Delpazolid, sutezolid, and telacebec (Q203)]
  Oxazolidinone and respiratory chain inhibitor candidates in Phase 2 evaluation. These target the cytochrome oxidase pathway recently characterized as the mechanism of pretomanid's bactericidal activity.

  \item[TBI-223 and TBAJ-587]
  Novel agents in early clinical development targeting different bacterial pathways to broaden the drug-resistant TB treatment repertoire.

\end{description}

% ============================================================
\section{Elimination Efforts}
\label{sec:elimination}
% ============================================================

\subsection{WHO End TB Strategy (2015--2035)}

The End TB Strategy, adopted by the World Health Assembly in 2014, sets milestone targets relative to a 2015 baseline:
\begin{itemize}
  \item \textbf{2020:} 20\% reduction in TB incidence; 35\% reduction in TB deaths.
  \item \textbf{2025:} 50\% reduction in TB incidence; 75\% reduction in TB deaths; 0\% of TB-affected families facing catastrophic costs.
  \item \textbf{2030:} 80\% reduction in incidence; 90\% reduction in deaths.
  \item \textbf{2035:} 90\% reduction in incidence; 95\% reduction in deaths; full elimination of catastrophic TB costs.
\end{itemize}
The 2020 milestones were missed, primarily due to COVID-19 disruption of health services and, in high-burden settings, pre-existing structural challenges.

\subsection{Global Plan to End TB 2023--2030}

The Stop TB Partnership's Global Plan provides costed targets and resource requirements. For 2023--2030, key targets include reaching 45 million people with TB treatment, 45 million with TPT, and achieving full rapid diagnostic coverage. Estimated annual funding required rises from approximately US\$22 billion by 2027 to US\$35 billion by 2030. Current funding levels remain far below these requirements.

\subsection{UN High-Level Meeting Commitments (2023)}

The 2023 UN Political Declaration on TB renewed commitments from the 2018 declaration, setting 2027 targets:
\begin{itemize}
  \item 90\% treatment coverage including 4.5 million children and 1.5 million DR-TB patients.
  \item 90\% TPT coverage for eligible household contacts and PLHIV.
  \item 100\% access to rapid molecular diagnostics.
  \item US\$22 billion/year for TB services; US\$5 billion/year for R\&D.
\end{itemize}

\subsection{National Elimination Programs}

Several countries have adopted TB elimination targets (defined as incidence $<$1/100,000):
\begin{itemize}
  \item \textbf{Low-incidence countries} (U.S., most of Western Europe, Australia) pursue elimination through systematic LTBI testing and treatment in high-risk groups, enhanced immigration screening, and contact investigation.
  \item \textbf{High-burden countries} (India, Indonesia, China, Nigeria, Pakistan, South Africa, etc.) focus on universal drug susceptibility testing, community-based active case-finding, and TPT scale-up, often through the USAID TIFA framework or Global Fund grants.
  \item \textbf{India's National TB Elimination Programme (NTEP)} aims for elimination by 2025 -- a target retained despite COVID-19 setbacks -- through the Ni-kshay patient support scheme, universal DST, and Pradhan Mantri TB Mukt Bharat Abhiyaan community engagement. Website: \url{https://tbcindia.gov.in}
\end{itemize}

% ============================================================
\section{Institutions Using Mathematical Models for TB}
\label{sec:modeling_groups}
% ============================================================

Mathematical modeling is integral to TB policy: it quantifies the latent reservoir, projects the impact of interventions, estimates $R_0$, and supports cost-effectiveness analyses submitted to funding bodies. Below is a curated list of active groups.

\begin{description}[style=nextline]

  \item[TB Modelling Group, LSHTM (London, UK)]
  Led by Prof.\ Richard White and Dr.\ Rein Houben. Uses ODE-based and Bayesian models to quantify the latent burden, TB spectrum, and intervention impact. Hosts the Secretariat of the TB Modelling and Analysis Consortium (TB-MAC). Developed the TIME Impact tool (in Spectrum) for national TB programs. Produced WHO policy briefs on BCG, new vaccines, and post-2015 targets. \\
  Website: \url{https://github.com/lshtm-tbmg}

  \item[TB Modeling \& Translational Epi Group, Johns Hopkins Bloomberg SPH (Baltimore, USA)]
  Led by Prof.\ David Dowdy. Combines mechanistic/ODE models with economic evaluation, field epidemiology, and implementation science. Prominent outputs include analyses of subclinical TB transmission, programmatic case-finding efficiency, and TB/HIV co-epidemic dynamics. \\
  Website: \url{https://modeltb.org}

  \item[MRC Centre for Global Infectious Disease Analysis, Imperial College London (London, UK)]
  Contributed multi-model analyses of WHO End TB targets (the 11-model study of South Africa, China, and India). Develops age-structured TB-HIV models calibrated to African countries. Broader disease modeling infrastructure shared with COVID-19, malaria, and other pathogens. \\
  Website: \url{https://www.imperial.ac.uk/mrc-global-infectious-disease-analysis}

  \item[Centre for Mathematical Modelling of Infectious Diseases (CMMID), LSHTM (London, UK)]
  TB is a major research theme within this broader centre. CMMID provides shared methodology (fitting frameworks, sensitivity analysis tools) used by the TB Modelling Group. \\
  Website: \url{https://www.lshtm.ac.uk/research/centres/centre-mathematical-modelling-infectious-diseases}

  \item[Harvard T.H.\ Chan School of Public Health -- TB Epi Group (Boston, USA)]
  Epidemiological and mathematical modeling of TB in high-burden countries, including agent-based models and network approaches. Work on TB diagnostics value-of-information and cost-effectiveness. \\
  Website: \url{https://www.hsph.harvard.edu/epidemiology/}

  \item[Institute for Health Metrics and Evaluation (IHME), University of Washington (Seattle, USA)]
  Produces Global Burden of Disease (GBD) estimates for TB incidence, prevalence, mortality, and risk factors using statistical models. The GBD TB envelope is widely used as an input to transmission models. \\
  Website: \url{https://www.healthdata.org}

  \item[University of Cape Town -- DSI-NRF Centre of Excellence in Epidemiological Modelling and Analysis (SACEMA) (Cape Town, South Africa)]
  Focuses on TB modeling in the context of high HIV burden in Sub-Saharan Africa. TB-HIV interaction models, age-structured analyses, and TPT impact projections informing South African National TB Programme. \\
  Website: \url{https://www.sacema.org}

  \item[Centre for the Mathematical Modelling of Infectious Diseases, Erasmus MC (Rotterdam, Netherlands)]
  European TB modelling contributions; cross-border transmission analysis, immigration screening modelling, and analyses supporting ECDC guidance. \\
  Website: \url{https://www.erasmusmc.nl}

  \item[TB Modelling and Analysis Consortium (TB-MAC)]
  Not a single institution but a global collaboration whose secretariat is hosted at LSHTM. Brings together modeling groups worldwide to coordinate methods, data sharing, and policy-relevant analyses. \\
  Website: \url{https://tb-mac.org}

  \item[WHO Global TB Programme -- analytical team]
  Uses transmission models internally and in commissioned analyses to project global and regional incidence trajectories, assess impact of the End TB Strategy, and produce the annual Global TB Report estimates. \\
  Website: \url{https://www.who.int/teams/global-tuberculosis-programme}

  \item[CDC Division of Tuberculosis Elimination -- TB Epidemiologic Studies Consortium (TBESC)]
  Conducts mathematical and statistical modelling for U.S. TB elimination analyses, evaluating LTBI testing and treatment programs and projecting time-to-elimination under various scenarios. \\
  Website: \url{https://www.cdc.gov/tb/}

\end{description}

\end{document}
